% Part:sets-functions-relations
% Chapter: size-of-sets
% Section: pairing

\documentclass[../../../include/open-logic-section]{subfiles}

\begin{document}

\olfileid[ar]{sfr}{siz}{pai}
\olsection{دوال المزاوجة والرموز}

\begin{explain}
تجعل طريقة كانتور المتعرجة قابلية $\Nat^n$ للتعداد ظاهرةً للعيان.
لكن لنركز على مصفوفتنا التي تمثل $\Nat^2$. فباتباع الخط المتعرج في
المصفوفة وعدّ المواضع، يمكننا التحقق من أن $\tuple{1,2}$ يقترن
بالعدد~$7$. غير أنه سيكون من المفيد أن نستطيع حساب ذلك بطريقة
أكثر مباشرة. أي سيكون من المفيد أن يكون في متناولنا \emph{معكوس}
تعداد المسار المتعرج، $g\colon \Nat^2 \to \Nat$، بحيث
\[
g(\tuple{0,0}) = 0, \;
g(\tuple{0,1}) = 1, \;
g(\tuple{1,0}) = 2, \; \dots,
g(\tuple{1,2}) = 7, \; \dots
\]
وسيتيح لنا ذلك أن نحسب بالضبط أين يظهر $\tuple{n, m}$ في تعدادنا.

والواقع أننا نستطيع تعريف $g$ مباشرةً بالاستناد إلى ملاحظتين. أولًا:
إذا كانت خانة الصف ذي الرتبة $n$ والعمود ذي الرتبة $m$ تحتوي
القيمة~$v$، فإن خانة الصف ذي الرتبة $(n+1)$ والعمود ذي الرتبة
$(m-1)$ تحتوي القيمة $v + 1$. ثانيًا: يتألف الصف الأول من تعدادنا
من الأعداد المثلثية، ابتداءً بـ$0$، $1$، $3$، $6$، وهكذا. والعدد
المثلثي ذو الرتبة $k$ هو مجموع الأعداد الطبيعية $\leq k$، ويمكن حسابه
بالصيغة $k(k+1)/2$. وبضم هاتين الملاحظتين، تأمل الدالة الآتية:
\[
  g(n,m) = \frac{(n+m+1)(n+m)}{2} + n
\]
\textbf{تنبيه تصحيحي على الأصل (C031-01).} الحد $k(k+1)/2$ يجمع الطبيعيّات التي
لا تتجاوز $k$، لا التي تقل عنه؛ فأصلحنا حد الجمع وأبقينا دالة المزاوجة.
غالبًا ما نكتب $g(n, m)$ بدلًا من $g(\tuple{n, m})$ فحسب، لأنه أيسر
على العين. ويطلب منك هذا أولًا تعيين العدد المثلثي ذي الرتبة
$(n+m)^\text{th}$، ثم إضافة $n$ إليه. وهو يملأ المصفوفة بالطريقة
التي نريدها تمامًا. ولذلك، على وجه الخصوص، يرسل الزوج
$\tuple{1, 2}$ إلى $\frac{4 \times 3}{2} +
1 = 7$.

هذه الدالة $g$ هي \emph{معكوس} تعداد مجموعة من الأزواج. وتسمى مثل
هذه الدوال \emph{دوال مزاوجة}.
\end{explain}

\begin{defn}[دالة المزاوجة]
  تكون الدالة $f\colon A \times B \to \Nat$ \emph{دالة مزاوجة}
  حسابية إذا كانت $f$ متباينة. ونقول أيضًا إن $f$ \emph{ترمّز}
  $A \times B$، وإن $f(x,y)$ هو \emph{رمز} $\tuple{x,y}$.
\end{defn}

\begin{explain}
يمكننا استعمال دوال المزاوجة لترميز أزواج الأعداد الطبيعية مثلًا؛
وبعبارة أخرى، يمكننا تمثيل كل \emph{زوج} من العناصر بعدد
\emph{واحد}. وباستعمال معكوس دالة المزاوجة، يمكننا \emph{فك ترميز}
العدد، أي معرفة الزوج الذي يمثله.
\end{explain}

\begin{prob}
أعط تعدادًا لمجموعة جميع الأعداد النسبية غير السالبة.
\end{prob}

\begin{prob}
أثبت أن $\Rat$ !!{enumerable}. وتذكر أنه يمكن كتابة كل عدد نسبي
في صورة كسر $z/m$، حيث $z \in \Int$ و$m \in \Nat^+$.
\end{prob}

\begin{prob}
عرّف تعدادًا لـ$\Bin^*$.
\end{prob}

\begin{prob}
تذكر من مقررك التمهيدي في المنطق أن كل جدول صدق ممكن يعبر عن دالة
صدق. وبعبارة أخرى، دوال الصدق هي جميع الدوال من $\Bin^k \to \Bin$
لعدد ما~$k$. أثبت أن مجموعة جميع دوال الصدق قابلة للتعداد.
\end{prob}

\begin{prob}
أثبت أن مجموعة جميع المجموعات الجزئية المنتهية لمجموعة اعتباطية لا
نهائية !!{enumerable} هي !!{enumerable}.
\end{prob}

\begin{prob}
يقال إن مجموعة جزئية $A$ من $\Nat$ \emph{متممتها منتهية} إذا وفقط إذا
وجدت مجموعة جزئية منتهية $F\subseteq\Nat$ بحيث
$A=\Nat\setminus F$؛ أي إذا وفقط إذا كانت $\Nat\setminus A$ منتهية.
ولتكن $I$ المجموعة
التي تكون !!{element}sها بالضبط المجموعات الجزئية المنتهية من
$\Nat$، والمجموعات الجزئية منها التي متمماتها منتهية. أثبت أن $I$
!!{enumerable}.
\end{prob}

\begin{editorial}
\textbf{تصويب في المصدر (OLSIZ-002).}
كانت العبارة الإنجليزية ناقصة الصلة النحوية؛ فصُرّح هنا بأن المتممة تؤخذ
في $\Nat$ لمجموعة جزئية منتهية منها، على ما تثبته الصيغة المكافئة التالية.
\end{editorial}

\begin{prob}
أثبت أن اتحاد عائلة !!{enumerable} من مجموعات !!{enumerable} يكون
مجموعةً !!{enumerable}. أي إنه متى كانت $A_1$، $A_2$، \dots{} مجموعات،
وكانت كل $A_i$ !!{enumerable}، فإن المجموعة
$\bigcup_{i=1}^\infty
A_i$، وهي اتحادها جميعًا، تكون أيضًا !!{enumerable}. [ملاحظة: هذا صعب!]
\end{prob}

\begin{prob}
لتكن $f \colon A \times B \to \Nat$ دالة مزاوجة اعتباطية. رتّب عناصر
$\ran{f}$ تصاعديًا وأعد فهرستها، ثم فك رموزها بمعكوس $f$ على مداه، وبرهن
أن ذلك يتيح بناء تعداد لـ$A \times B$ إذا كان حاصل الضرب غير خال. وبيّن
حكم الحالة الخالية على حدة؛ فإن كان الحاصل منتهيًا غير خال كرر آخر عنصر
عند اعتماد تعريف التعداد بدالة شاملة من $\PosInt$، أو افهرس عناصره بمقطع
ابتدائي عند اعتماد التعريف البديل.
\end{prob}

\textbf{تنبيه تصحيحي على الأصل (C031-02).} ليس معكوس دالة المزاوجة، من حيث هو،
تعدادًا؛ لأن مجاله $\ran{f}$ قد لا يكون $\Nat$ ولا $\PosInt$ ولا مقطعًا
ابتدائيًا. ولذلك صُرّح بإعادة فهرسة الرموز وبالحالتين المنتهية والخالية.

\begin{prob}
حدد دالة ترمّز $\Nat^3$.
\end{prob}

\end{document}
